\section{Native Branch-and-Cut Solver}
\label{sec:native-branch-and-cut}

This section documents the toolbox's native branch-and-cut solver from three angles at once: mathematical model, algorithmic workflow, and code-level realization. The implementation is centered on
\texttt{include/hacdcpf/engine/branch_and_cut.hpp},
\texttt{src/engine/branch_and_cut.cpp},
\texttt{src/engine/bc\_utils.cpp},
\texttt{src/engine/bc\_branching.cpp},
\texttt{src/engine/bc\_cuts.cpp}, and
\texttt{src/engine/bc\_relaxation.cpp}, while the LP kernel relies on
\texttt{include/hacdcpf/engine/dual_simplex.hpp} and
\texttt{src/engine/dual_simplex.cpp}.
The public namespace is \texttt{hacdcpf::engine};
\texttt{include/hacdcpf/solver/branch_and_cut.hpp} re-exports the same solver for the adapter layer.

The discussion below is intentionally code-aligned. When the implementation differs from a textbook branch-and-cut design, this section follows the current C++ code rather than the idealized label.

\subsection{Scope and Code Map}

The module currently serves two problem classes:
\begin{itemize}
  \item MILP: the full branch-and-cut path with root cuts, pseudo-cost / reliability branching, cut pool, solution pool, primal heuristics, sequential tree search, and heterogeneous parallel tree search.
  \item MINLP: a lighter spatial branch-and-bound path with NLP relaxations, branching, and pruning, but without the MILP path's cut generation, cut pool, or dedicated cut worker.
\end{itemize}

The main source files are summarized in Table~\ref{tab:bc-file-map}.

\begin{table}[H]
\centering
\small
\caption{Branch-and-cut code map}
\label{tab:bc-file-map}
\begin{tabularx}{\linewidth}{L{0.32\linewidth}L{0.62\linewidth}}
\toprule
File & Main responsibility \\
\midrule
\texttt{include/hacdcpf/engine/branch_and_cut.hpp} & Public API, \texttt{BCOptions}, \texttt{BCStats}, \texttt{BCResult}, and solver enums. \\
\texttt{include/hacdcpf/engine/detail/bc\_types.hpp} & Core data structures: nodes, pseudo-costs, cut requests, LP relaxation records, and thread-side payloads. \\
\texttt{include/hacdcpf/engine/detail/bc\_pools.hpp} & Global \texttt{CutPool} and \texttt{SolutionPool}. \\
\texttt{include/hacdcpf/engine/detail/bc\_threading.hpp} & Thread-safe queues, shared incumbent, shared cut/solution pools, and atomic parallel statistics. \\
\texttt{src/engine/bc\_utils.cpp} & Objective normalization, fractional-variable detection, MILP presolve, node bound propagation, reduced-cost fixing, and row insertion. \\
\texttt{src/engine/bc\_branching.cpp} & Most-infeasible branching, pseudo-cost branching, and reliability branching probes. \\
\texttt{src/engine/bc\_cuts.cpp} & Tableau GMI/Gomory cuts, cover cuts, legacy integer-rounding cuts, and cut dispatch. \\
\texttt{src/engine/bc\_relaxation.cpp} & LP and NLP relaxation solves plus auxiliary rounding logic. \\
\texttt{src/engine/branch_and_cut.cpp} & Root strengthening, sequential B\&C, heterogeneous parallel B\&C, and MILP/MINLP top-level orchestration. \\
\texttt{src/engine/dual_simplex.cpp} & Standard-form construction, sparse and dense simplex, dual-simplex reoptimization, and incremental bound updates in standard form. \\
\bottomrule
\end{tabularx}
\end{table}

The most direct solver entry points in the current toolbox are:
\begin{itemize}
  \item \texttt{solve\_milp\_bc(prob, opt)} for MILP.
  \item \texttt{solve\_minlp\_bc(prob, opt)} for MINLP.
  \item \texttt{NativeBranchAndCutAdapter} for integration into the solver-adapter interface.
  \item \texttt{analysis/topology\_analysis.cpp} for ONR-type MILP usage.
  \item \texttt{analysis/time\_series\_pf.cpp}, \texttt{tests/test\_unit\_commitment.cpp}, and \texttt{tests/test\_uc\_performance.cpp} for UC and scheduling-style MILP integration and tuning.
\end{itemize}

\subsection{Active and Prepared Tree Data Structures}

The active MILP and MINLP solve paths still use the legacy \texttt{Node} structure:
\begin{itemize}
  \item full node bounds \texttt{lb}/\texttt{ub},
  \item relaxation solution \texttt{x\_relax},
  \item optional seed \texttt{x\_seed},
  \item simplex warm-start hint \texttt{basis\_hint},
  \item node bound \texttt{bound}, and
  \item tree depth.
\end{itemize}

A recent implementation change made \texttt{Node::branch\_child()} explicitly lightweight: it copies only bounds, basis hint, bound, and depth, but skips copying \texttt{x\_relax} and \texttt{x\_seed}. This matters because branching creates children frequently, and the solver now avoids unnecessary vector copies on the hot path.

At the same time, \texttt{bc\_types.hpp} now also contains a prepared diff-based representation:
\begin{itemize}
  \item \texttt{BoundChange},
  \item \texttt{TreeNode}, and
  \item \texttt{TreeDomainManager}.
\end{itemize}

These structures support LCA-based switching between nodes by replaying bound changes rather than copying full bound vectors. As of the current implementation, they are present in the type layer but are not yet wired into the active branch-and-cut search loop. Therefore the operational solver still uses the legacy full-bound \texttt{Node} path.

\subsection{Problem Class and Outer Mathematical Model}

For MILP, the input object is \texttt{MIPModel}, whose linear part is \texttt{LPModel}:
\begin{align}
  \min/\max \quad & c^\top x \\
  \text{s.t.}\quad & A x \le b, \\
                    & A_{eq} x = b_{eq}, \\
                    & \ell \le x \le u, \\
                    & x_j \in \mathbb{Z} \text{ or } \{0,1\}, \quad j \in \mathcal{I}.
\end{align}

The outer B\&C logic compares all bounds in a unified minimization metric. The helper
\texttt{objective\_value(c, x, sense)} returns \(c^\top x\) for minimization and
\(-c^\top x\) for maximization, so pruning, pseudo-cost updates, and gap computation can all be written as
\begin{align}
  z_{LP}(N) &\le z^\star \le z_{inc}, \\
  \text{gap} &= \max\!\left(0, \frac{z_{inc} - z_{LB}}{\max(1, |z_{inc}|)}\right).
\end{align}

For a fractional integer variable \(x_j^\ast\), the two children are
\begin{align}
  \text{down child}:\quad & x_j \le \lfloor x_j^\ast \rfloor, \\
  \text{up child}:\quad   & x_j \ge \lceil x_j^\ast \rceil.
\end{align}

A node is pruned if any of the following happens:
\begin{itemize}
  \item the node bounds are inconsistent, that is, \(\ell_i > u_i\) for some \(i\);
  \item the LP or NLP relaxation is infeasible or numerically rejected;
  \item the node bound is not better than the incumbent, that is, \(z_{LP}(N) \ge z_{inc}\);
  \item the relaxation solution is already integral, in which case it becomes a candidate incumbent immediately.
\end{itemize}

\subsection{MILP Presolve and Bound Strengthening}

\texttt{milp\_presolve()} runs before the root node is solved. It is not a general-purpose black-box presolve; it is a branch-and-cut-oriented propagation engine focused on bounds, fixings, and matrix simplification.

\paragraph{1. Inequality and equality bound tightening}
For an inequality row \(a^\top x \le b\), the code computes the minimum activity
\begin{equation}
L(a,\ell,u)=\sum_{a_j>0} a_j \ell_j + \sum_{a_j<0} a_j u_j,
\end{equation}
and then derives tighter upper or lower bounds variable by variable. Equality rows use both minimum and maximum activity, so they can be stronger. Integer and binary bounds are rounded immediately with \(\lfloor \cdot \rfloor\) or \(\lceil \cdot \rceil\).

\paragraph{2. Singleton-row processing}
If a row contains only one nonzero coefficient, it directly implies a bound on that variable. Equality singletons can fix a variable outright. Fixed singleton rows are then removed from the active matrix.

\paragraph{3. Redundant-row elimination}
If an inequality already satisfies \(\max a^\top x \le b\) under current bounds, the row is redundant and removed from \texttt{A}.

\paragraph{4. Binary probing inside presolve}
For unfixed binary variables, the presolver probes \(x_j=0\) and \(x_j=1\), applies three rounds of propagation in each branch, and then:
\begin{itemize}
  \item fixes the variable if one side is infeasible, or
  \item merges implied bounds back into the root model if both sides are feasible.
\end{itemize}

\paragraph{5. Fixed-variable substitution}
If \(\ell_j=u_j\), the contribution of that variable is absorbed into \(b\) and \(b_{eq}\), and the corresponding matrix column is zeroed out during sparse-matrix rebuild. This reduces effective LP density before simplex begins.

A lighter version of the same idea is repeated at tree nodes through
\texttt{node\_bound\_propagation()}. The node-level propagator does not rebuild matrices; it only tightens current node bounds for two or three rounds before the LP solve, which allows cheap infeasibility detection before simplex reoptimization.

\subsection{Standard Form and LP Relaxation Kernel}

LP node solves are not performed directly on
\((A,b,A_{eq},b_{eq},\ell,u)\). The solver first constructs a bounded standard form via \texttt{build\_standard\_form\_lp()}.
For each original variable, it defines
\begin{equation}
  x = \ell + y, \qquad y \ge 0,
\end{equation}
and stores finite upper bounds in \texttt{var\_ub}, that is,
\(0 \le y_j \le u_j-\ell_j\).
The resulting internal LP is represented as a maximization problem,
\begin{align}
  \max \quad & c_{max}^\top z \\
  \text{s.t.}\quad & \widetilde A z = \widetilde b, \\
                    & 0 \le z \le \widetilde u.
\end{align}

Important implementation details are:
\begin{itemize}
  \item The internal objective is always maximization, with \(c_{max}=-c_{min}\).
  \item \(\le\) rows receive slack variables; \(\ge\) rows receive surplus plus artificial variables; equality rows receive artificial variables.
  \item If a shifted right-hand side is negative, the whole row is multiplied by \(-1\), and the row sense is flipped so that the stored RHS becomes nonnegative.
  \item Finite upper bounds are kept implicitly through \texttt{var\_ub} and \texttt{at\_upper}; they are not converted into explicit rows.
\end{itemize}

So the LP kernel is closer to a bounded-form revised simplex than to the textbook lower-bound-only standard form. This is the key reason why both bounded-form GMI generation and dual-simplex warm starts are efficient in the current implementation.

\paragraph{Fast bound updates}
Most children modify only one or two bounds. For that reason,
\texttt{update\_standard\_form\_bounds()} updates only
\texttt{lb\_shift}, \texttt{b}, \texttt{var\_ub}, and \texttt{objective\_const}, without rebuilding \texttt{A}. When no more than 50 lower-bound components change, the code applies sparse column-wise incremental updates; otherwise it falls back to a full sparse matrix-vector product.

\paragraph{LP relaxation solve modes}
\texttt{solve\_lp\_relaxation()} has two front-end modes:
\begin{itemize}
  \item Default path: \texttt{use\_simplex\_lp\_nodes=true}. The solver uses the native simplex kernel and returns a basis hint, reduced costs, and optionally a persisted sparse basis object for child-node warm starts and GMI separation.
  \item Fallback path: \texttt{use\_simplex\_lp\_nodes=false}. The LP is wrapped as an NLP with zero Hessian and solved by the native IPM. This still gives a relaxation solution, but tableau-based cuts are unavailable unless a \texttt{SimplexResult} exists.
\end{itemize}

\subsection{Warm Start by Dual Simplex}

The main theoretical fact behind the tree search is simple: when branching only changes variable bounds, the constraint matrix \(A\) remains unchanged. Therefore the parent basis is often still structurally useful for the child. In bounded-form simplex, reduced costs are determined by
\begin{equation}
  \bar c = c_{max} - A^\top y, \qquad y = B^{-T} c_B,
\end{equation}
so a basis reused under changed bounds often stays dual feasible while losing primal feasibility. That is exactly the setting in which dual simplex is effective.

The implementation currently contains three practical warm-start regimes.

\paragraph{Path A: sparse warm start for large LPs}
For \(m > \texttt{kSparseSimplexThreshold} = 200\), the code prefers a sparse basis path built around \texttt{SparseBasis}. The main steps are:
\begin{itemize}
  \item reuse the parent's persisted sparse basis and reduced costs when available;
  \item validate basis indices and block artificial columns from entering;
  \item pre-clean the basis by swapping out basic artificial variables with row slacks whenever possible, which avoids false infeasibility after bound changes;
  \item recompute \(x_B\), reduced costs, objective value, and \texttt{at\_upper} status with \texttt{sparse\_full\_recompute()} or the equivalent sparse operations;
  \item if primal feasibility is violated but dual feasibility is retained, run \texttt{sparse\_dual\_simplex\_reoptimize()};
  \item if dual reoptimization fails numerically but the basis can be refactorized, fall back to sparse primal Phase II as a last warm-start rescue.
\end{itemize}

When sparse warm start succeeds, the solver persists the sparse basis again for descendants instead of computing a dense inverse.

\paragraph{Path B: dense warm start for small LPs}
For \(m \le 200\), the solver uses the dense inverse path. If a cached inverse is available, it reconstructs the basis state quickly; otherwise it recomputes the state from the basis itself. If the reconstructed state is primal infeasible but dual feasible, it calls \texttt{dual\_simplex\_reoptimize()}.

\paragraph{Path C: cold start when warm start is unavailable or rejected}
The simplex kernel still supports cold-start Phase I / Phase II solves with slack and artificial variables. For large LPs, the sparse cold-start path uses sparse primal simplex and may use a Big-M objective if crash replacement does not remove all artificials. For small LPs, the dense path uses a classic Phase I / Phase II sequence.

\paragraph{A deliberate engineering tradeoff at tree nodes}
In the main MILP tree search, \texttt{SimplexOptions.allow\_cold\_start} is set to \texttt{false}. Therefore, when a large child-node warm start fails, the code does \emph{not} attempt an expensive cold-start Phase I / Phase II solve. Instead, that node LP is marked as failed and the branch-and-cut layer prunes it.
This is a speed-first engineering choice: it reduces node cost substantially on large UC-style MILPs, but strict completeness depends on warm-start robustness. Root solves and root cut re-solves still retain cold-start capability so that the algorithm does not lose the model at the very beginning.

\paragraph{Incumbent-aware LP early cutoff}
Both sequential and parallel MILP paths provide an incumbent bound to the simplex solver. The LP kernel can then terminate early when the node's LP bound is already provably worse than the incumbent, which avoids many unnecessary pivot iterations in nodes that would be pruned anyway.

\subsection{Cut Families and Admission Rules}

The current cut module implements three distinct families.

\paragraph{1. Tableau-based Gomory / GMI cuts}
\texttt{add\_basis\_gomory\_cuts()} builds bounded-form GMI cuts from the current simplex basis. Conceptually, if row \(r\) corresponds to a basic integer variable with value \(\beta\) and fractional part
\(f_0 = \beta - \lfloor \beta \rfloor\), the code first extracts the tableau row
\begin{equation}
  \rho = e_r^\top B^{-1} A,
\end{equation}
using either a dense \(B^{-1}\) row or sparse BTRAN. It then computes bounded-form GMI coefficients for each nonbasic column, taking into account:
\begin{itemize}
  \item whether the original variable is integer or continuous,
  \item whether that variable is currently represented at its upper bound, and
  \item the shifted lower-bound representation used by the standard form.
\end{itemize}

After that, slack and surplus columns are eliminated algebraically, and the cut is mapped back to the original variable space as
\begin{equation}
  a_{cut}^\top x \le b_{cut}.
\end{equation}

Candidate GMI cuts are not inserted blindly. They are filtered by:
\begin{itemize}
  \item violation at the current LP point,
  \item efficacy \(\text{viol}/\lVert a_{cut} \rVert_2\),
  \item density cap \texttt{gmi\_max\_density},
  \item a binary-activity score that favors cuts acting on highly fractional binary support,
  \item a binary-support ratio filter, and
  \item near-parallel rejection against already selected cuts through \texttt{gmi\_max\_parallelism}.
\end{itemize}

The ranking score is
\begin{equation}
  \text{score} = \text{efficacy} \cdot \left(1 + w_{act} \cdot \text{activity}\right),
\end{equation}
where \(w_{act}\) is \texttt{gmi\_activity\_weight}.

\paragraph{2. Cover cuts}
\texttt{add\_cover\_cuts()} scans inequality rows that contain only positive coefficients on binary variables. For a cover set \(C\) with
\begin{equation}
  \sum_{j\in C} a_j > b,
\end{equation}
it generates the classical cover inequality
\begin{equation}
  \sum_{j\in C} x_j \le |C| - 1.
\end{equation}
The implementation builds the cover greedily by sorting coefficients in descending order and accumulating until the right-hand side is exceeded.

\paragraph{3. Legacy row-based integer-rounding cuts}
\texttt{add\_mir\_like\_cuts()} reads fractional right-hand sides directly from original inequality rows and builds an older integer-rounding cut from row coefficients. In the current public interface, this family is exposed primarily by \texttt{CutType::IntRounding}.

\paragraph{Actual mapping of \texttt{CutType}}
Table~\ref{tab:bc-cut-map} records the \emph{actual current code behavior}, not the ideal meaning suggested by enum names.

\begin{table}[H]
\centering
\small
\caption{Cut-type enum and current implementation mapping}
\label{tab:bc-cut-map}
\begin{tabularx}{\linewidth}{L{0.18\linewidth}L{0.28\linewidth}L{0.44\linewidth}}
\toprule
Enum value & Actual dispatch target & Notes \\
\midrule
\texttt{None} & No cuts & Baseline mode. \\
\texttt{IntRounding} & \texttt{add\_mir\_like\_cuts} & Legacy row-based integer-rounding cuts on original inequalities. \\
\texttt{MIR} & \texttt{add\_basis\_mir\_cuts} & \texttt{add\_basis\_mir\_cuts()} currently delegates directly to \texttt{add\_basis\_gomory\_cuts()}. \\
\texttt{Gomory} & \texttt{add\_basis\_gomory\_cuts} & Main tableau-based GMI/Gomory generator. \\
\texttt{Cover} & \texttt{add\_cover\_cuts} & Binary knapsack-style cover cuts. \\
\texttt{All} & Gomory $\rightarrow$ Cover & The dispatcher now skips the MIR alias in this mode to avoid duplicate tableau GMI passes. \\
\bottomrule
\end{tabularx}
\end{table}

So if an upper-layer module says it is using ``MIR cuts'' through \texttt{CutType::MIR}, the current implementation should be understood as using the tableau GMI/Gomory generator under a different enum label. The only truly row-based family currently exposed is \texttt{IntRounding}.

\subsection{Cut Pool and Solution Pool}

\paragraph{CutPool}
The global \texttt{CutPool} stores cuts that were already generated and can be cheaply re-separated later in the tree. Each pool entry keeps a dense coefficient vector, right-hand side, age, best efficacy placeholder, and norm.
Its current policy is:
\begin{itemize}
  \item duplicate rejection by cosine similarity above 0.95,
  \item age-based replacement when the pool is full, and
  \item age reset for violated cuts plus purge after \texttt{cut\_pool\_max\_age} consecutive non-violations.
\end{itemize}

\paragraph{SolutionPool}
The solution pool stores the best \(K\) feasible integer solutions and supports guided rounding. For a fractional integer variable, guided rounding uses a vote across stored solutions to choose the rounding direction. This is more structure-aware than plain \(\text{round}(x_j)\) because it reuses historically successful integer patterns.

\subsection{Branching, Node Selection, and Primal Heuristics}

\paragraph{Branch-variable selection}
The code currently supports three branching strategies:
\begin{itemize}
  \item \texttt{FirstFractional}: branch on the first fractional variable. Useful as a baseline or debugging strategy.
  \item \texttt{MostInfeasible}: pick the variable with fractional part closest to 0.5. The score is
        \(0.5 - |\text{frac}(x_j)-0.5|\).
  \item \texttt{Pseudocost}: use a product score
        \begin{equation}
          s_j = \max\!\left(f_j \bar q_j^{\downarrow}, \varepsilon\right)
                \max\!\left((1-f_j) \bar q_j^{\uparrow}, \varepsilon\right),
          \qquad f_j = x_j - \lfloor x_j \rfloor,
        \end{equation}
        where \(\bar q_j^{\downarrow}\) and \(\bar q_j^{\uparrow}\) are average pseudo-cost estimates.
\end{itemize}

When pseudo-cost branching is enabled, shallow nodes may use
\texttt{choose\_branch\_var\_reliability()} first. That routine probes unreliable candidates by solving temporary down/up child LPs, updates pseudo-costs with true gains, and only then selects the branch variable. The current implementation behaves as follows:
\begin{itemize}
  \item reliability branching is used only up to depth 10;
  \item candidates that already received root probing data require fewer additional observations;
  \item the probing budget is more aggressive at the root and lighter at deeper nodes.
\end{itemize}

\paragraph{Node selection}
The public API exposes \texttt{BestFirst}, \texttt{DepthFirst}, and \texttt{Hybrid}. The active implementation of \texttt{Hybrid} is slightly more specific than the enum name suggests:
\begin{itemize}
  \item before an incumbent exists, new nodes are pushed to DFS;
  \item after an incumbent exists, new nodes are pushed to the best-bound priority queue;
  \item however, if the DFS stack is nonempty, it is still popped before the priority queue.
\end{itemize}
So the current \texttt{Hybrid} strategy is best read as DFS-biased with best-bound backfill, not as a strict one-time DFS-to-best-first mode switch.

\paragraph{Primal heuristics used in the main solve path}
The main branch-and-cut driver currently employs the following primal heuristics:
\begin{itemize}
  \item quick rounding after the root LP,
  \item user warm-start acceptance and warm-start repair,
  \item a fix-and-solve style feasibility pump at the root,
  \item LP diving at the root,
  \item direct rounding and guided rounding at child nodes,
  \item crossover from the two best solution-pool incumbents, and
  \item a simple RINS-like heuristic in the sequential path.
\end{itemize}

The helper \texttt{try\_rounding\_heuristic()} still exists in \texttt{bc\_relaxation.cpp}, but the production MILP path in \texttt{branch\_and\_cut.cpp} now performs the main rounding logic inline.

\subsection{Root-Node Strengthening Pipeline}

In the current implementation, the root node is much more than ``solve one LP and start branching''. The actual pipeline is shown in Algorithm~BC-R1.

\paragraph{Algorithm BC-R1: Root strengthening}
\begin{algorithmic}[1]
\Procedure{RootStrengthen}{$prob, opt$}
  \State run \texttt{milp\_presolve} on the \texttt{LPModel}
  \State solve the root LP relaxation to obtain \(x_0^{LP}\), a basis hint, and the root bound
  \If{cuts are enabled}
    \For{$k=1$ to \texttt{root\_cut\_rounds}}
      \State generate up to \texttt{cuts\_per\_round} root cuts
      \If{no cut was added} \State \textbf{break} \EndIf
      \State extend the basis hint to accommodate new slack/artificial columns
      \State re-solve the cut-augmented root LP
      \If{the re-solve fails} \State roll back the added root cuts and stop root cutting \EndIf
      \If{bound improvement stalls} \State stop root cutting early \EndIf
    \EndFor
  \EndIf
  \State seed the cut pool from accepted root cuts
  \State try user warm start and warm-start repair
  \State run the feasibility-pump loop
  \State run LP diving
  \State run two-phase root probing: domain propagation, then LP probing
  \State if an incumbent exists, run reduced-cost fixing at the root
  \State push the strengthened root node into the search queue
\EndProcedure
\end{algorithmic}

Several implementation details are especially important.

\paragraph{1. Root cut re-solve basis extension}
After adding root cuts, the standard-form dimensions change. The old basis hint is therefore not directly compatible. The code explicitly rebuilds the basis index array so that:
\begin{itemize}
  \item old inequality rows keep their relative order,
  \item new cut rows are inserted between the old inequality block and the equality block,
  \item old equality rows are shifted downward, and
  \item slack, surplus, and artificial column indices are renumbered consistently.
\end{itemize}
This is one of the key robustness mechanisms behind successful root cut re-solves.

\paragraph{2. Root cut rollback and sparse-basis cleanup}
If a root cut round adds cuts but the re-solve fails, the code truncates the newly appended inequality rows and restores the pre-cut basis, root relaxation point, and root bound. Even after successful root cutting, sparse-basis pointers created during cut re-solves are cleared before the fresh base standard form is built, because they may otherwise reference already-destroyed local matrices.

\paragraph{3. Stall detection}
Root cutting stops early if the relative bound improvement falls below \(10^{-4}\) for two consecutive rounds.

\paragraph{4. Warm-start repair}
If \texttt{prob.initial\_solution} exists, the solver first clamps it to bounds and rounds integer variables. If it is not feasible, the solver can still attempt repair by fixing the integer variables to the warm-start pattern and re-solving an LP over the continuous variables.

\paragraph{5. Feasibility pump}
The current root feasibility-pump logic has two stages:
\begin{itemize}
  \item a cheap direct rounding test, and
  \item up to five iterations of ``round integer variables, fix them, re-solve the restricted LP for the continuous part''.
\end{itemize}
So the current implementation is closer to a fix-and-solve repair loop than to the more elaborate distance-minimization feasibility-pump variants found in the literature.

\paragraph{6. LP diving}
The root diving heuristic starts from the root LP solution and greedily fixes near-integral variables while repeatedly re-solving warm-started LPs. It:
\begin{itemize}
  \item is skipped when a warm start already provided an incumbent,
  \item scales its LP budget with model size,
  \item uses a time budget of at most roughly 3\% of the global time limit and no more than about 1 second,
  \item batch-fixes variables that are already very close to integrality, and
  \item tries the opposite rounding direction if the nearest-integer fix fails.
\end{itemize}

\paragraph{7. Two-phase root probing}
Root probing has two distinct phases.
\begin{itemize}
  \item Domain-propagation probing: for each eligible unfixed binary variable, test \(x_j=0\) and \(x_j=1\) by bound propagation only. If one side becomes inconsistent, fix the variable to the other side.
  \item LP probing: for the most fractional unfixed binaries, solve temporary down/up LPs. This both initializes pseudo-costs and may produce additional fixings when one side is LP-infeasible or already dominated by the incumbent.
\end{itemize}

The LP probing budget is adaptive. It is smaller for large models and smaller again when a good incumbent already exists.

\paragraph{8. Reduced-cost fixing}
If an incumbent exists, the root can use reduced costs to fix nonbasic integer variables. The current implementation computes
\begin{equation}
  \Delta = z_{inc} - z_{LP}(N),
\end{equation}
and for a nonbasic integer variable at one of its bounds, fixes it at that bound if
\begin{equation}
  |rc_j| \ge \Delta.
\end{equation}
In the present code, reduced costs are handled in the internal minimization metric and applied only to nonbasic integer variables that are already sitting at their lower or upper bound.

\subsection{Child Processing and Sequential Tree Search}

For each child node, the helper \texttt{process\_single\_child()} applies the following steps in order:
\begin{enumerate}
  \item check bound consistency;
  \item run \texttt{node\_bound\_propagation()} on the node bounds;
  \item update the standard-form bounds by \texttt{update\_standard\_form\_bounds()} without rebuilding the matrix;
  \item solve the child LP with the parent basis hint;
  \item if an incumbent exists, run reduced-cost fixing and re-check consistency;
  \item if the LP has at most 500 inequality rows, separate violated cuts from the global cut pool and re-solve the child LP with those pool cuts;
  \item if \texttt{cuts != None} and the child depth is within \texttt{max\_cut\_depth}, generate new node cuts, add them to the global cut pool, and re-solve the node on a transient cut-augmented LP;
  \item attempt direct rounding and guided rounding to search for new incumbents.
\end{enumerate}

The 500-row threshold is a pure engineering switch. Above that size, copying \texttt{LPModel}, rebuilding \texttt{StandardFormLP}, and re-solving the LP for pool separation becomes too expensive relative to the expected gain, so the implementation skips pool separation.

A second important detail is that local node cuts are handled on a transient cut-augmented LP copy. If that re-solve succeeds, the node keeps the improved bound and LP point, but the persistent basis hint is restored to the pre-cut basis. This design keeps descendants aligned with the base LP plus node bounds rather than with a private cut-augmented submodel.

The sequential main loop can be summarized as Algorithm~BC-S1.

\paragraph{Algorithm BC-S1: Sequential MILP B\&C}
\begin{algorithmic}[1]
\Procedure{SequentialBC}{$root$}
  \While{time limit not reached, node limit not reached, and queue not empty}
    \State pop the next live node \(N\)
    \If{$z_{LP}(N) \ge z_{inc}$} \State \textbf{continue} \EndIf
    \If{$x^{LP}(N)$ is integral}
      \State update incumbent and solution pool
      \State \textbf{continue}
    \EndIf
    \State choose branch variable by reliability branching or pseudo-cost branching
    \State create down and up children
    \State process both children with \texttt{process\_single\_child}
    \State push the more promising child to DFS first and the other to the global queue structure
    \State periodically run crossover and the RINS-like heuristic
    \State update live-node lower bound and global gap
    \If{relative gap is below tolerance} \State terminate \EndIf
  \EndWhile
\EndProcedure
\end{algorithmic}

In the sequential path:
\begin{itemize}
  \item crossover is attempted every 100 explored nodes;
  \item the RINS-like heuristic is attempted every 200 explored nodes;
  \item the simplex kernel receives an atomic incumbent bound so that poor nodes can be cut off early during LP reoptimization itself.
\end{itemize}

\subsection{Heterogeneous Parallel Tree Search}

The parallel MILP path is not a homogeneous ``all workers do the same thing'' implementation. It is intentionally heterogeneous:
\begin{itemize}
  \item if \texttt{opt.cuts != None} and \texttt{max\_cut\_depth > 0}, the solver runs one dedicated cut worker plus \(N-1\) explorer threads;
  \item if cuts are only allowed at the root, all threads become explorers because a cut worker would be idle.
\end{itemize}

The explorer threads are split into two roles:
\begin{itemize}
  \item \texttt{Diver}: incumbent-oriented, DFS-biased, pseudo-cost branching, and local plunging.
  \item \texttt{Prover}: bound-oriented, best-bound queue preference, reliability branching at shallow depths, and stronger focus on proving optimality.
\end{itemize}

\paragraph{Shared parallel structures}
The parallel code uses the following shared objects:
\begin{itemize}
  \item \texttt{ThreadSafeNodeQueue}: a shared queue containing both a best-bound priority queue and a DFS stack.
  \item \texttt{SharedIncumbent}: objective bound shared by atomic compare-and-swap and solution vector shared by mutex.
  \item \texttt{SharedCutPool}: a \texttt{shared\_mutex}-protected cut pool.
  \item \texttt{SharedSolutionPool}: a mutex-protected solution pool.
  \item \texttt{CutRequestQueue}: a bounded explorer-to-cut-worker request queue.
  \item \texttt{AtomicBCStats}: lock-free or low-lock parallel counters.
\end{itemize}

\paragraph{Explorer behavior}
Each explorer owns a private copy of the base standard form. This avoids repeated synchronization around LP matrix data. In addition:
\begin{itemize}
  \item a Diver maintains a thread-local DFS stack;
  \item the maximum size of that local DFS plunge is now controlled by
        \(\texttt{kLocalDfsMax} = \max(1, \texttt{opt.max\_plunge\_depth})\);
  \item when the local plunge stack is full, additional children are spilled to the shared queue;
  \item explorers still perform node LP solves, reduced-cost fixing, pool separation, rounding, and guided rounding locally.
\end{itemize}

So \texttt{max\_plunge\_depth} is no longer a dormant option: it now actively controls how aggressively Diver threads keep nodes in thread-local DFS before returning them to the shared queue.

\paragraph{Cut worker behavior}
Explorer threads do not perform expensive tableau separation themselves in the parallel path. Instead they package the current \texttt{SimplexResult}, node bounds, node depth, and LP relaxation point into a \texttt{CutRequest} and submit it non-blockingly to the cut worker.
If the request queue is full, the request is dropped. This affects cut opportunities, not solution correctness.

The cut worker then:
\begin{itemize}
  \item pops cut requests,
  \item generates cuts,
  \item inserts newly generated rows into the shared cut pool, and
  \item when idle, continues to search for incumbents via guided rounding and crossover from the shared solution pool.
\end{itemize}

\paragraph{A code-level nuance: parallel tree cutting is currently Gomory-only}
There is an important implementation asymmetry here. Root cutting and sequential child cutting both use the generic \texttt{add\_cuts()} dispatcher, so they honor the configured \texttt{CutType}. The dedicated parallel cut worker, however, currently calls \texttt{add\_basis\_gomory\_cuts()} directly. Therefore tree-level cut generation in the parallel path is effectively Gomory/GMI-only, even if the user selected \texttt{MIR}, \texttt{Cover}, or \texttt{All}. This is a real code fact and should be kept in mind when interpreting parallel runs.

\paragraph{Monitor loop}
The main thread acts as a monitor. Roughly every millisecond it checks:
\begin{itemize}
  \item wall-clock time limit,
  \item node limit,
  \item optimality gap, and
  \item whether the shared queue is empty and no explorer is still active.
\end{itemize}

The last condition is verified twice with a short sleep in between, which reduces false termination due to queue/worker races.

\subsection{MINLP Path}

\texttt{solve\_minlp\_bc()} is currently a spatial branch-and-bound rather than a full branch-and-cut engine:
\begin{itemize}
  \item the root and child relaxations are solved by \texttt{solve\_nlp\_relaxation()} through the native IPM,
  \item branching still reuses the same branch-variable selection framework and pseudo-cost table,
  \item there is no cut generation, no cut pool, no solution pool, and no heterogeneous parallel explorer/cut-worker design.
\end{itemize}

So the name ``branch-and-cut'' is fully accurate for the MILP path, but the MINLP path should currently be understood as branch-and-bound with NLP relaxations.

\subsection{Options and Statistics}

Table~\ref{tab:bc-options} summarizes the most important user-visible solver options.

\begin{table}[H]
\centering
\small
\caption{Key \texttt{BCOptions} fields}
\label{tab:bc-options}
\begin{tabularx}{\linewidth}{L{0.27\linewidth}L{0.15\linewidth}X}
\toprule
Option & Default & Purpose \\
\midrule
\texttt{max\_nodes} & 50000 & Maximum number of explored tree nodes. \\
\texttt{max\_lp\_iter} & 500 & LP/NLP iteration budget parameter used in relaxation solves. \\
\texttt{time\_limit\_sec} & 120.0 & Global wall-clock time limit. \\
\texttt{int\_tol} & \(10^{-5}\) & Integrality tolerance. \\
\texttt{gap\_tol} & \(10^{-4}\) & Relative optimality-gap tolerance. \\
\texttt{lp\_tol} & \(10^{-6}\) & LP tolerance used to derive simplex feasibility and optimality tolerances. \\
\texttt{root\_cut\_rounds} & 5 & Maximum number of root cut rounds. \\
\texttt{cuts\_per\_round} & 10 & Maximum new cuts per round or per node. \\
\texttt{max\_cut\_depth} & 0 & Tree nodes deeper than this do not generate fresh cuts. \\
\texttt{cuts} & \texttt{None} & Cut-family selection switch. \\
\texttt{branching} & \texttt{Pseudocost} & Branch-variable selection strategy. \\
\texttt{node\_sel} & \texttt{Hybrid} & Node-selection strategy. \\
\texttt{use\_simplex\_lp\_nodes} & true & Whether node LPs use the simplex path with basis reuse. \\
\texttt{use\_feasibility\_pump} & true & Whether the root feasibility-pump loop is enabled. \\
\texttt{cut\_pool\_max\_size} & 2000 & Maximum cut-pool capacity. \\
\texttt{cut\_pool\_max\_age} & 50 & Age threshold for purging stale pool cuts. \\
\texttt{solution\_pool\_size} & 10 & Number of best incumbents kept for guided rounding and crossover. \\
\texttt{num\_threads} & 1 & Thread count; values above 1 enable the heterogeneous parallel path. \\
\texttt{cut\_worker\_queue\_size} & 64 & Maximum number of pending explorer-to-cut-worker requests. \\
\texttt{max\_plunge\_depth} & 10 & Active limit on Diver thread-local DFS plunging before spillback to the shared queue. \\
\texttt{gmi\_*} & several thresholds & Controls GMI efficacy, density, near-parallel rejection, binary activity, and support filters. \\
\bottomrule
\end{tabularx}
\end{table}

\texttt{BCStats} records diagnostics about the actual search behavior, as summarized in Table~\ref{tab:bc-stats}.

\begin{table}[H]
\centering
\small
\caption{Meaning of \texttt{BCStats} fields}
\label{tab:bc-stats}
\begin{tabularx}{\linewidth}{L{0.27\linewidth}X}
\toprule
Field & Meaning \\
\midrule
\texttt{nodes\_explored} & Number of explored nodes. \\
\texttt{lp\_solves} & Total number of LP or NLP relaxation solves. \\
\texttt{cuts\_added} & Number of newly generated cuts inserted into models. \\
\texttt{best\_bound} & Current global lower bound over live nodes. \\
\texttt{best\_obj} & Current incumbent objective value. \\
\texttt{gap} & Relative MIP gap in the solver's internal minimization metric. \\
\texttt{runtime\_sec} & Total runtime. \\
\texttt{pool\_cuts\_separated} & Number of cuts reintroduced from the global cut pool. \\
\texttt{crossover\_incumbents} & Number of incumbents found by the crossover heuristic. \\
\texttt{cut\_requests\_submitted} & Number of cut requests sent from explorers to the cut worker. \\
\texttt{cut\_requests\_dropped} & Number of cut requests dropped because the queue was full. \\
\texttt{cut\_worker\_cuts} & Number of cuts generated by the dedicated cut worker. \\
\texttt{cut\_worker\_incumbents} & Number of incumbents found by the cut worker through guided rounding or crossover. \\
\texttt{status} & Solver status string, aligned with \texttt{SolveStats.status}. \\
\bottomrule
\end{tabularx}
\end{table}

\subsection{Applications, Tests, and Code-Aligned Caveats}

\paragraph{1. Current in-toolbox usage}
The native MILP branch-and-cut solver is currently used most directly in:
\begin{itemize}
  \item ONR and topology-reconfiguration workflows,
  \item UC and scheduling-style MILP integration, and
  \item performance and regression tests that compare the native implementation against external solvers.
\end{itemize}

The current UC-oriented tuning used by \texttt{analysis/time\_series\_pf.cpp} is representative:
\begin{itemize}
  \item \texttt{cuts = CutType::Gomory},
  \item \texttt{root\_cut\_rounds = 15},
  \item \texttt{cuts\_per\_round = 20},
  \item \texttt{use\_feasibility\_pump = false},
  \item \texttt{gap\_tol = 1e-3},
  \item large \texttt{max\_lp\_iter}, and
  \item parallelism up to 8 hardware threads when available.
\end{itemize}

\paragraph{2. Existing test coverage}
The current implementation is directly exercised by at least the following test groups:
\begin{itemize}
  \item \texttt{tests/test\_branch\_and\_cut.cpp} for correctness on small MILP/MINLP instances,
  \item \texttt{tests/test\_unit\_commitment.cpp} for UC-style integration and solver behavior,
  \item \texttt{tests/test\_uc\_performance.cpp} for larger UC scaling and parallel performance, and
  \item the root-cut re-solve robustness tests associated with basis extension behavior.
\end{itemize}

\paragraph{3. Important caveats for users and developers}
The current codebase contains several behaviorally important caveats.
\begin{itemize}
  \item \textbf{\texttt{CutType::MIR} is not a distinct MIR implementation.} It currently aliases the tableau GMI/Gomory generator.
  \item \textbf{\texttt{CutType::All} no longer duplicates MIR/GMI dispatch.} It now executes one tableau GMI pass plus cover cuts.
  \item \textbf{Large child-node LPs usually do not cold start.} This is a deliberate speed optimization, not a purely textbook completeness-first design.
  \item \textbf{Pool separation is intentionally skipped for larger LPs.} Once the LP is large, a second full re-solve per child becomes too expensive.
  \item \textbf{Parallel cut requests may be dropped.} This weakens cut generation opportunities but does not invalidate incumbents or bounds already proven.
  \item \textbf{The dedicated parallel cut worker is currently Gomory-only.} Parallel tree cuts therefore do not perfectly mirror the configured \texttt{CutType} semantics.
  \item \textbf{Diff-based tree management is prepared but not active.} \texttt{TreeNode} and \texttt{TreeDomainManager} exist in the type layer, but the running solver still uses the legacy full-bound \texttt{Node} representation.
\end{itemize}

\paragraph{4. Practical interpretation}
Taken together, the native branch-and-cut solver is not a minimal teaching implementation. It is an engineering-oriented MILP core that combines:
\begin{itemize}
  \item bounded-form simplex with dual-simplex warm reoptimization,
  \item an aggressively strengthened root node,
  \item several incumbent-finding heuristics,
  \item reusable cut and solution pools,
  \item sequential and heterogeneous parallel tree search, and
  \item application-driven tuning for power-system optimization models such as ONR and UC.
\end{itemize}

The most natural next algorithmic extensions, if stricter generality is desired, are a truly distinct MIR implementation, a more explicit recovery path when large-node warm starts fail, activation of diff-based tree-state management, and richer cut or parallel logic for the MINLP path.
